\scanpage{47}

II\hfill\mspage{- 11 -}\hfill n$^\circ$338

A tout $A$-module $M$, projectif de rang 1, faisons correspondre un élément
$J$ de $\underline I$ isomorphe à $M$ (prop.5), et notons $\operatorname{Cl}(M)$ l'image canonique de $J$ dans
$\underline I/\underline P$ ; comme deux idéaux fractionnaires isomorphes (en tant que $A$-modules)
sont homothétiques, $\operatorname{Cl}(M)$ est indépendant du choix de $J$. Pour que deux
$A$-modules $M$ et $M'$ projectifs, de type fini, et de rang 1 soient \uline{isomorphes},
il faut et il suffit que les éléments $\operatorname{Cl}(M)$ et $\operatorname{Cl}(M')$ soient \uline{égaux}. Comme
le produit tensoriel de deux modules projectifs est un modules projectif, et
donc sans torsion, il résulte de ce qui a été vu au début du n$^\circ$ que la multiplication
dans le groupe $\underline I/\underline P$ correspond au \uline{produit tensoriel} des $A$-modules
projectifs de rang 1. On dit parfois que $\underline I/\underline P$ est le groupe des $A$-modules
projectifs de rang 1.

\begin{center}\texttt{-:-:-:-:-}\end{center}

\margin{46}
